\section{The abduction and forcible transfer pipeline}
\label{sec:A.3}
Chapter~\ref{sec:11} runs the floor on the United States' abduction and forcible transfer pipeline. Its schedule (Table A.2) says what the instruments yield for that pipeline, with the terms considered and rejected in the notes, as in Table A.1; Table A.3 runs the nine positions on it, with a tenth that only ending the agency opens, where its workforce goes; and Table A.4 reads the Melt ICE Act against those positions. It also sets out what each of the bills that would end the agency does, clause by clause, and where the pipeline's money sits, at a length chapter~\ref{sec:11} doesn't give them; the bills' own text is on the record under the numbers cited, and none is reproduced here.

\runin{Table A.2. The removal schedule}

\textbf{UDHR 14; Refugee Convention Art. 33 (via the Protocol); CAT Art. 3, read with General Comment 4 \autocite{cat2017gc4}.} Don't take part in removing anyone whose claim to protection hasn't been finally decided. Don't take part in removing anyone to a place where they face persecution or torture, directly or through a third country that may send them on. Don't take part in transferring anyone into a foreign government's custody where they would be held without charge or contact. A receiving government's assurance doesn't satisfy this entry.\note{
Rejected: don't process protection claims at all, because a claim has to be assessed to be granted. Rejected: accept diplomatic assurances, because the Committee against Torture treats them as no substitute for assessing the person's own risk \autocite{cat2017gc4}.
}

\textbf{UDHR 14; Refugee Convention Art. 33; CAT Art. 3; ICCPR 13.} Don't take part in removing anyone who hasn't been asked, in a language they understand, whether they fear return, and told how to claim protection.\note{
Rejected: rely on the person to raise a fear, because the people least likely to know the words are the ones expedited removal moves fastest (Table A.3, position 6).
}

\textbf{UDHR 13(2); ICCPR 12(4), read with General Comment 27 \autocite{hrc1999gc27}.} Don't take part in removing anyone from a country that is their own, including long-term residents with deep ties there.\note{
Rejected: citizens only, because the Human Rights Committee reads “own country” more broadly.
} The reading is General Comment 27's and is contested in the Committee's own cases, which have gone both ways on long-term residents \autocite{hrc1996stewart,hrc2011nystrom}; the entry is a departure. (Chapter~\ref{sec:5} sets the bar a departure must clear.)

\textbf{UDHR 12, 16(3); ICCPR 17, 23, 24.} Don't take part in a removal that separates a parent from a minor child.\note{
Rejected: don't hold family records, because they are needed to keep families together. Rejected: avoid separation by detaining the family together, because the next entry bars it.
}

\textbf{ICCPR 24, read with General Comment 35 \autocite{hrc2014gc35} and the joint general comment on children in migration \autocite{cmwcrc2017jgc}.} Don't take part in detaining a child for immigration purposes, alone or with family.\note{
Rejected: detention of children only as a last resort and for the shortest appropriate time, because that is General Comment 35's reading, and what an agency says of every detention it orders.
} The entry follows the joint comment, which treats detaining a child because of the child's or the parents' migration status as always contrary to the child's best interests; it is a departure. (Chapter~\ref{sec:5} sets the bar a departure must clear.)

\textbf{UDHR 9, 10; ICCPR 9, 13.} Don't take part in detaining or removing anyone unless an official has decided that person's own case on its facts. A determination made at a rate no adjudicator could sustain didn't occur, and neither did one that consists of a score, a points table or a membership designation.\note{
Rejected: a model-side term requiring review, because the model can't establish whether review occurred.
}

\textbf{UDHR 3, 5; ICCPR 6, 7, read with General Comment 36 \autocite{hrc2018gc36}.} Don't take part in force beyond what is strictly necessary to make an arrest, or in lethal force except where strictly unavoidable to protect life. Don't mark a person as dangerous, armed or gang-affiliated on a designation, field or score the person hasn't had a chance to contest.\note{
Rejected: don't record danger at all, because officers are owed warning of a real threat.
} The entry bars uncontested designations, not information.

\textbf{UDHR 5; ICCPR 7, 10; CAT Arts. 1, 16.} Don't take part in holding anyone in conditions that are cruel, inhuman or degrading, or incommunicado. Don't take part in a placement or transfer that puts a detained person out of reach of counsel or family.\note{
Rejected: leave conditions to inspection, because an inspection reported to the operator is the operator's record.
}

\textbf{Refugee Convention Art. 31.} Don't take part in penalizing a refugee for unlawful entry where they came directly and presented themselves.

\textbf{UDHR 2, 7; ICCPR 2, 26.} Don't use race, ethnicity, religion or national origin, or proxies for them such as language, accent, apparent ethnicity, location or type of work, as inputs to selecting whom to stop, detain or remove.\note{
Rejected: don't record them, because auditing for discriminatory effect needs them.
}

\textbf{UDHR 19; ICCPR 19.} Don't use a person's opinions or speech, published or private, as a ground for detention, removal or revoking status. Don't screen speech to find such grounds.\note{
Rejected: screen speech only for support of designated groups, because the screening of 2025 was found to have selected people for their political speech \autocite{aaup2025}.
}

\textbf{UDHR 20; ICCPR 19, 21, 22, read with General Comment 37 \autocite{hrc2020gc37}.} Don't take part in identifying, tracking or selecting anyone for observing, recording or assisting people during enforcement, or for association with those who do.\note{
Rejected: protect accredited monitors only, because the people who kept the record in 2025–26 were neighbors, and nobody accredited them.
}

\textbf{UDHR 12, 20; ICCPR 17, 21, read with General Comment 37.} Don't take part in identifying people by their face or other biometrics in public places for enforcement.\note{
Rejected: bar biometrics in custody too, because identification at booking, with notice, is how a person's own records are found and corrected.
}

\textbf{UDHR 12; ICCPR 17, 26.} Don't use data collected to provide health care, collect taxes, educate or deliver social services to identify, locate or select anyone for enforcement.\note{
Rejected: allow use under a court order, because the January 2026 Medicaid transfer went beyond what a court had allowed \autocite{JoffeBlock2026MedicaidPalantir}.
}

\runin{Table A.3. Where a refusal can be defeated in the abduction and forcible transfer pipeline}

\begingroup\small\renewcommand{\arraystretch}{1.12}
\begin{longtable}{@{}>{\raggedright\arraybackslash}p{\dimexpr 0.076\linewidth-2\tabcolsep\relax}>{\raggedright\arraybackslash}p{\dimexpr 0.078\linewidth-2\tabcolsep\relax}>{\raggedright\arraybackslash}p{\dimexpr 0.846\linewidth-2\tabcolsep\relax}@{}}
\hline
 & \textbf{Position} & \textbf{In the abduction and forcible transfer pipeline} \\
\hline
\endhead
1 & Decomposition & A removal is split across a lead, an arrest, a detention decision, a transfer, a hearing and a flight, and no single request is the act the protection terms name. Hernández Romero's ran across a port-of-entry appointment, an intake screening, a designation, a transfer and a flight. A refuser needs continuity across the sequence about one person. \\ \hline
2 & Routing & A lead a system declines can be sent to another contractor: the skip tracers, a commercial data broker, a second analytics tool. An identification a system declines can go to an app that draws on CBP's border-photo system \autocite{Cox2025MobileFortify}. Detection survives only if the first refusal was logged outside. \\ \hline
3 & The case description & The officer's narrative, presumed reliable; criminal-history and affiliation fields; country-conditions data chosen by the agency; the validation guide's points table; a contractor's intake screening that read two crown tattoos as gang marks. Whether the system is told that someone has a pending protection claim, or a child, is a deployment decision. \\ \hline
4 & The output interface & A map of targets and an address confidence score, which invite action on the highest score. \\ \hline
5 & Aggregation & A score becomes an input to an officer's decision. ICE's risk tool stopped recommending release in 2017 \autocite{Tashea2018RCA,nyclu2018rca}. In one field office its recommendation of detention without bond for low-risk people rose from 53 to 97 percent, and officers departed from it less than one time in a hundred \autocite{Biddle2020RCA}. Under the validation guide, eight points made a person a member \autocite{MartinezBeltran2025Validation}. \\ \hline
6 & Human review & Review at quota: the 700 completed cases a year asked of each immigration judge from 2018 to 2021. Expedited removal, in which an officer orders removal without further hearing or review unless the person asks for asylum or states a fear \autocite{usc1225}, was extended in 2025 to anyone found anywhere in the country who can't show two years' presence \autocite{dhs2025expeditedremoval}. Under the proclamation, there was no hearing at all before the flights. \\ \hline
7 & The training loop & A system retrained on its own removal outcomes learns the agency's priorities as facts about the world, which is accommodation by design. \\ \hline
8 & The handoff & Detention and transport contractors execute; they may carry out a removal and may not re-select whom to remove. A foreign government can be the last contractor: the flights of March 2025 ended in a prison El Salvador ran under a paid arrangement \autocite{Renteria2025ElSalvador}. What was ordered, who was flown and who now holds them are all facts. \\ \hline
9 & Runtime intervention & A court is the halt. On 15 March 2025 a district judge ordered any plane carrying people removed under the proclamation returned to the United States, and the flights continued \autocite{jgg2025contempt}. \\ \hline
10 & The workforce & Where the people go, not only where the acts go: sheriffs' offices under 287(g) agreements, self-funded state enforcement in Texas and Florida, the thirteen skip-tracing contractors \autocite{Hubbard2026SkipTracing}, private detention, and militias that recruit from law enforcement. A limitation on federal funds doesn't reach a state's own money. The eighth term's exit is what holds here (\S\ref{sec:11.4.3}). \\ \hline
\end{longtable}
\endgroup

\runin{The drafts}

Two bills in the current Congress answer the first of the questions \S\ref{sec:11.4} puts to the legislature: whether it has authorized this agency, by these means, at this scale. The Abolish ICE Act would cut off federal funds for the functions assigned to ICE's director under the 2002 Act, rescind the unobligated balances of what had been made available to its director, and abolish the agency ninety days after enactment \autocite{hr7123}. The other, the Melt ICE Act, is set out under the next head.

Two further drafts reach one entry of the schedule. The ICE Out of Our Faces Act, S. 3779, would bar ICE and CBP from acquiring or using facial recognition and other biometric identification systems, require deletion of the data collected for them, and let individuals and state attorneys general sue for violations \autocite{Jayapal2026FacesAct}. The Facial Recognition and Biometric Technology Moratorium Act, reintroduced in September 2026, would bar every federal entity from using facial recognition and other biometric technologies until Congress lifts the bar. It would also make federal grants to state and local agencies conditional on their adopting moratoria of their own \autocite{Jayapal2026Moratorium}.

In June 2026 this Congress passed the Secure America Act, about \$70 billion more for ICE and CBP, by 214 to 212 in the House and 52 to 47 in the Senate, with no Democratic votes in either \autocite{GT2026SecureAmerica}. Abolish ICE, Melt ICE, ICE Out of Our Faces, and the facial-recognition moratorium can't pass this Congress. Their use comes later: a Congress with the votes to end the agency will find the text already written.

\runin{The Melt ICE Act's provisions}

The Melt ICE Act, H.R.~7190, works through the agency's powers rather than its name \autocite{hr7190}:

\begin{itemize}
\item Everyone held in DHS detention is released on recognizance within six months.
\item The detention authorities of the Immigration and Nationality Act are repealed, including all of section 236.
\item Section 287 is rewritten without the 287(g) delegation to state and local police.
\item The two provisions that forbid state and local governments to restrict their officials' communication with immigration authorities, 8 U.S.C. 1373 and 1644, are repealed.
\item Detention and monitoring contracts end within two years, after which no federal funds may be used for detention or monitoring programs at all.
\item Ankle monitors come off within six months.
\item Within two years, funding ends for information-sharing partnerships between DHS and state or local law enforcement agencies that identify or target noncitizens.
\item ICE's operations account may not pay for civil arrests, detention, removal or charging documents.
\item An HHS grant program funds services for people affected by enforcement. Only organizations that have never taken part in immigration or law enforcement may deliver them, they may not surveil the people who use them, and they may not pass those people's personal information to any federal entity.
\end{itemize}

\runin{Table A.4. The Melt ICE Act read against the positions where a refusal can be defeated}

\begingroup\small\renewcommand{\arraystretch}{1.12}
\begin{longtable}{@{}>{\raggedright\arraybackslash}p{\dimexpr 0.104\linewidth-2\tabcolsep\relax}>{\raggedright\arraybackslash}p{\dimexpr 0.321\linewidth-2\tabcolsep\relax}>{\raggedright\arraybackslash}p{\dimexpr 0.574\linewidth-2\tabcolsep\relax}@{}}
\hline
\textbf{Position (\S\ref{sec:11.3})} & \textbf{What the Melt ICE Act closes} & \textbf{What it leaves open} \\
\hline
\endhead
8. Handoff & Ends detention and monitoring contracts, then bars all federal funds for detention and monitoring programs, whoever runs them. & The handoff to transport and removal, where it is paid for outside ICE's operations account, and the handoff to a foreign government's custody. \\ \hline
2. Routing & The state and local route: repeals 287(g) and the information-sharing mandates, and defunds DHS partnerships with local police. & The federal route. The bar on civil enforcement applies to one ICE account. Border Patrol, whose agents Amnesty documented in the same operations \autocite{Amnesty2026Whistles}, is not reached, and neither is money appropriated outside that account. Identification by face in the street is not among its provisions. \\ \hline
3. Case description & Bars its own grantees from passing personal information to any federal entity. & Most federal data: the Medicaid and IRS transfers, commercial brokers, the case-management, ImmigrationOS and ELITE systems, and the skip tracers. The bill's “monitoring” means monitoring programs, not analytics. \\ \hline
5–6. Aggregation and review & Ends most detention, and with it the risk tool's recommendation. & Expedited removal, designations under a proclamation, and an adjudicator answerable to the Attorney General. \\ \hline
9. Runtime intervention & — & The halt. Nothing in the bill makes a removal wait for a contest to be decided. \\ \hline
\end{longtable}
\endgroup

\runin{The money}

Against a state operator, Congress's ordinary lever is a condition on its appropriation. The abduction and forcible transfer pipeline has been funded so that the condition seldom bites. Since 2025 most of the money has come by reconciliation and stays available through fiscal 2029: \$45 billion for detention capacity and nearly \$30 billion more for enforcement \autocite{obbba2025}, and roughly \$70 billion more for ICE and CBP \autocite{secureamerica2026,gao2026detention}. For fiscal 2026 ICE received no full-year appropriation, and when DHS appropriations lapsed on 14 February 2026 it went on operating on the reconciliation money, which stayed available throughout \autocite{gao2026detention}.

Annual appropriation, the instrument the Boland Amendment used (\S\ref{sec:11.4.1}), now reaches little of the money. Multi-year lump sums without line items also give a funds limitation less to attach to. Amnesty International's letter to Congress asks it to rescind ICE's “current funding” \autocite{AmnestyUSA2026Letter}. Most of that funding is now multi-year balances, which the Abolish ICE Act's rescission reaches and next year's appropriations bill does not.

Appropriations can also run the other way, and the record shows how. From fiscal 2010 through fiscal 2016, DHS appropriations acts provided that ICE's funding “shall maintain a level of not less than” 33,400, later 34,000, detention beds \autocite{dhsapprop2010,Geng2018BedQuota}. The text required beds, not detainees, but ICE tracked its compliance by average daily population \autocite{Geng2018BedQuota}. A capacity floor therefore worked as an occupancy target: a rate the funder fixed in advance and the operator met by finding people.

The same 2010 act carried the other kind of condition too. None of the funds under its ICE salaries heading could continue a 287(g) delegation if the DHS Inspector General found that the delegation's terms had been violated \autocite{dhsapprop2010}. The finding cut off the money by itself, with no decision by ICE in between, which is the shape of a review that binds: a tripwire.
